\documentclass[../main.tex]{subfiles}

\begin{document}
    \noindent \hspace*{0.75cm}Berdasarkan hasil praktikum pada Modul {\thechapter} tentang \leftmark, dapat diambil kesimpulan bahwa:
    \begin{enumerate}
        \item {
            {\Flowchart} merupakan sebuah diagram yang digunakan untuk menggambarkan alur kerja sebuah proses atau algoritma. {\Flowchart} dapat memudahkan pemahaman sebuah algoritma.
        }
        \item {
            {\Pseudocode} merupakan sebuah bahasa yang bisa digunakan untuk menjelaskan sebuah program kepada orang yang tidak memiliki pengetahuan atau tidak mengerti bahasa pemrograman. Penulisan {\pseudocode} dapat ditulis dalam format apapun asal bisa dimengerti oleh pembaca.
        }
        \item {
            Bahasa pemrograman C++ merupakan bahasa pemrograman tingkat tinggi yang berarti C++ merupakan bahasa pemrograman yang dekat dengan bahasa manusia. C++ memiliki beberapa {\header} {\file} bawaan yang memudahkan penulisan kode seperti ``{\crnfamily iostream}'' yang menyediakan objek-objek dan fungsi-fungsi yang dibutuhkan untuk berintraksi dengan {\ioi} dan {\ioo} {\stream}.
        }
    \end{enumerate}
\end{document}
